\documentclass[10pt]{article}

\usepackage{setspace}
\usepackage{amsmath, amssymb, graphicx}
\usepackage{physics}
\usepackage{booktabs, array}
\usepackage[a4paper, total={16cm, 24.5cm}]{geometry}
\usepackage{cancel}
\usepackage[english]{babel}
\usepackage[utf8]{inputenc}
\usepackage{graphicx}
\usepackage{float}
\usepackage{enumitem}
\usepackage{xfrac}
\usepackage{listings, hyperref}
\usepackage{color}
\usepackage{multirow}
\usepackage{longtable}
\usepackage{bm}

\newcommand{\ie}{{\it i.e.},\ }
\def\eg{\emph{e.g.}\ }
\def\ea{\emph{et al.}}
\def\cf{\emph{c.f.}\ }
\newcommand{\pd}{{\phantom{\dagger}}}
\newcommand{\vk}{\mathbf{k}}
\newcommand{\vd}{\textbf{\textit{d}}}
\newcommand{\bsigma}{\boldsymbol{\sigma}}
\newcommand{\bc}{\mathbf{c}}
\newcommand{\ba}{\boldsymbol{a}}
\newcommand{\bs}{\mathbf{s}}
\newcommand{\bdelta}{\boldsymbol{\delta}}
\newcommand{\vdel}{\bm{\delta}}
\newcommand{\vc}{\bm{c}}
\newcommand{\bK}{\bm{K}}
\newcommand{\bR}{\bm{R}}
\newcommand{\bk}{\boldsymbol{k}}
\newcommand{\bq}{\boldsymbol{q}}
\newcommand{\br}{\boldsymbol{r}}
\newcommand\numberthis{\addtocounter{equation}{1}\tag{\theequation}}
\graphicspath{{/}{images/}}

\newcommand{\todo}[1]{\textcolor{red}{TODO: #1}}


\begin{document}

\title{QPI SC symmetry}
\date{\today}

\maketitle

\tableofcontents

\section{SC Basics}

BdG Hamiltonian:
\begin{align}
    H (\bk)  &= \begin{pmatrix}
        \varepsilon (\bk) & \Delta (\bk) \\
        \Delta^{*} (\bk) & -\varepsilon (\bk)
    \end{pmatrix} \\
    &= \begin{pmatrix}
        \frac{\Delta}{\sqrt{2E (E + \varepsilon )}} & \frac{\Delta}{\sqrt{2E (E - \varepsilon)}} \\
        \frac{-(E+\varepsilon)}{\sqrt{2E (E + \varepsilon )}} & \frac{E - \varepsilon}{\sqrt{2E (E - \varepsilon)}}
    \end{pmatrix}
    \begin{pmatrix}
        - E & 0 \\
        0 & E
    \end{pmatrix}
    \begin{pmatrix}
        \frac{\Delta^*}{\sqrt{2E (E + \varepsilon )}} &  \frac{-(E+\varepsilon)}{\sqrt{2E (E + \varepsilon )}}\\
        \frac{\Delta^*}{\sqrt{2E (E - \varepsilon)}} & \frac{E - \varepsilon}{\sqrt{2E (E - \varepsilon)}}
    \end{pmatrix} \\
    &= U (\bk) \Lambda (\bk) U^\dagger (\bk)
\end{align}
with $E = \sqrt{\varepsilon^2 + |\Delta|^2}$

Retarded Green's function:
\begin{align}
    G_0 (\omega, \bk) &= \frac{1}{ \omega + i \Gamma - H (\bk) } \\
    &= U (\bk) \frac{1}{ \omega + i \Gamma - \Lambda (\bk) } U^\dagger (\bk)
\end{align}

\section{QPI Basics}

We consider a local impurity in the presence of the superconducting phase defined by the Hamiltonian
\begin{equation}
    H = H_0 + V \, ,
\end{equation}
where $H_0$ is the unperturbed (impurity-free) superconducting Hamiltonian. $V$ is the local impurity term, given as
\begin{equation} \label{eqn:impurity_ham}
    V = \sum_{i \, \sigma \sigma'} [ \hat V ]_{\sigma \sigma'}^{\phantom{\dagger}} \delta^{\phantom{\dagger}}_{i, \, \boldsymbol{R} = 0} c^\dagger_{i \sigma} c^{\phantom{\dagger}}_{i \sigma'} \, ,
\end{equation}
where $\hat V = V_0 \, \hat I$ for a scalar impurity and $\hat V = V_0 \, \hat \sigma_z$ for a magnetic impurity%\,\cite{balatsky2006}.
The QPI spectra are calculated using the $T$-matrix formalism%\,\cite{hirschfeld2015, bunney2026},
\begin{equation} \label{eqn:qpi_rs}
    \delta \rho (\omega, \br) = - \frac{1}{\pi}  \text{Im} [  \text{Tr} [ \hat G_0 (\omega, \br) \, \hat T (\omega) \, \hat G_0 ( \omega, - \br) ] ] \, ,
\end{equation}
where the $T$-matrix can be shown to take the form:
\begin{equation}
    \hat T ( \omega ) = \frac{\hat V}{1- \hat G_0 (\omega, \br = 0) \hat V} \; .
\end{equation}
While this definition of the $T$-matrix is a formal definition, in practice it can be calculated as the non-trivial solution to the matrix equation
\begin{equation}
    (1- \hat G_0 (\omega, \br = 0) \hat V) \, \hat T ( \omega ) = \hat V\ .
\end{equation}
Note that the Born approximation $\hat T = \hat V$ corresponds to a first-order approximation of the binomial expansion of the $T$-matrix.
Subgap states (that exist in the superconducting gap) are associated with the poles of the $T$-matrix (at zero temperature)%\,\cite{callaway2013}.
The condition for the existence of these bound states is
\begin{equation} \label{eqn:bound_evals}
    \text{det} [ 1 - \hat G_0 (\omega, \br = 0) \, \hat V ] = 0 \, .
\end{equation}
Assuming a given unperturbed Green's function (Hamiltonian) and $\omega$, this fixes the value of $V_0$ in Eqn.\,\eqref{eqn:impurity_ham}. Note that in principle this could introduce large $V_0$ values, so accuracy could potentially be improved by including bands higher in energy now accessible to the impurity states.

\section{Minimal model}

We consider a triangular lattice, with hopping $t$ between nearest neighbour sites, and an onsite chemical potential $\mu$. We add a variety of nearest neighbour superconducting pairings, corresponding to the different irreps of $C_{6v}$ (which can be represented by Cooper pairs on the nearest neighbour bonds). For the moment, we will only consider even parity pairings, \ie the $A_1, A_2, E_2$ irreps.
In reciprocal space, the Hamiltonian is given associated
\begin{align}
    H &= \sum_{\bk, \sigma \sigma'}
    ( \varepsilon (\bk) - \mu ) \delta^{\phantom{\dagger}}_{\sigma \sigma'} c^{\dagger}_{\bk \sigma} c^{\phantom{\dagger}}_{\bk \sigma'}
    + \Delta (\bk) (i \sigma_y)^{\phantom{\dagger}}_{\sigma \sigma'} c^{\dagger}_{-\bk \sigma} c^{{\dagger}}_{\bk \sigma'}
    + \Delta^{*} (\bk) (-i \sigma_y)^{\phantom{\dagger}}_{\sigma \sigma'} c^{\phantom{\dagger}}_{-\bk \sigma} c^{\phantom{\dagger}}_{\bk \sigma'} \\
    \varepsilon (\bk) &= 2t \sum_{\alpha = 1}^{3} \cos (\ba_\alpha \cdot \bk ) \\
    \ba_1 &= (1,0), \hspace{1cm} \ba_2 = (-1/2, \sqrt{3}/2) \hspace{1cm} \ba_3 = -\ba_1 - \ba_2
\end{align}
The even parity superconducting pairings on nearest neighbours can be expressed as:
\begin{align}
    \Delta (\bk) &= 2 \sum_{\alpha = 1}^3 \Delta_\alpha \cos (\ba_\alpha \cdot \bk) \\
    \Delta^{A_1} &= (1,1,1)  &&(\text{extended } s)\\
    \Delta^{E_2, 1} &= (1,-1/2, -1/2) &&(d_{x^2-y^2}) \\
    \Delta^{E_2, 2} &= (0,-\sqrt{3}/2,\sqrt{3}/2)  &&(d_{xy})\\
    \Delta^{E_2, +} &= (1, e^{-2 \pi i /3}, e^{2 \pi i /3}) &&(d+id)\\
    \Delta^{E_2, -} &= (1, e^{2 \pi i /3}, e^{-2 \pi i /3}) &&(d-id)
\end{align}


\end{document}